\documentclass{situationist-preprint}

\title{Naturalized Epistemology and Normativity}
\msid{EPI-26-a}
\msfamily{EPI-26}
\msdiscipline{Epistemology}
\mstype{Principal manuscript}
\mskeywords{naturalized epistemology; Quine; normativity; reliability; ought implies can; cognitive psychology; justification}

\begin{document}
\maketitle

\begin{abstract}
Quine proposed that epistemology be replaced by psychology: the study of how evidence is related to theory should become the empirical study of how the human subject, given sensory input, produces a theory of the world. Kim objected that this abandons the normative question of what we ought to believe, and so changes the subject. We argue for a position that accepts half of each view. Epistemology cannot do without psychology, for two reasons that we make precise: the epistemic norms that bind human believers are constrained by what human believers can do, and the conditions of knowledge include facts about the reliability and safety of belief-forming processes that only empirical inquiry can establish. But the normative questions of epistemology are not psychological questions. No set of purely descriptive premises about how beliefs are formed settles which beliefs are justified without a normative premise linking the descriptive facts to epistemic standing, and that premise is not itself a finding of psychology. Epistemology is therefore continuous with psychology in its premises and its methods, while its conclusions remain normative.
\end{abstract}

\section{The problem}

\citet{Quine1969} argued that the foundationalist programme of justifying our knowledge of the external world from sense experience had failed, both in its doctrinal part, the deduction of truths about bodies from truths about experience, and in its conceptual part, the translation of talk of bodies into talk of experience. He concluded that we should give up the attempt to vindicate science from a standpoint outside it, and study instead, as part of science, how human beings arrive at science from the meagre input of their sensory surfaces. Epistemology, on this proposal, \enquote{simply falls into place as a chapter of psychology} \citep{Quine1969}.

\citet{Kim1988} replied that epistemology is a normative discipline: its central concept, justification, is evaluative, and to replace the question which beliefs are justified by the question which beliefs are caused by which stimulations is to stop doing epistemology. The causal question is legitimate, but it is a different question. Quine's own later statements moved toward a compromise: normative epistemology is the technology of truth-seeking, and its norms are hypothetical imperatives whose ends are truth and prediction and whose means are discovered empirically \citep{Quine1986}.

Between replacement and autonomy, several positions have been defended. \citet{Goldman1986} argued that epistemology has two components, a normative framework specifying the criteria of justification, which is the business of philosophy, and the application of those criteria to human cognition, which requires cognitive science. \citet{Kornblith2002} argued that knowledge is a natural phenomenon, studied by cognitive ethology, and that conceptual analysis has little role. \citet{BishopTrout2005} argued that epistemology should take its norms from the findings of the psychology of reasoning, in particular from the success of simple statistical prediction rules over expert judgment. \citet{Stich1990} argued that once the empirical facts about cognition are in view, the value of true belief itself becomes a matter of what we happen to want.

We argue that Goldman's division is correct in outline and that it can be made sharper. Section~2 sets out the framework, Section~3 proves that epistemology depends on psychology in two distinct ways, and Section~4 that its normative questions are not thereby psychological.

\section{Framework}

\begin{definition}[Descriptive and normative claims]\label{d:desc}
A \emph{descriptive epistemic claim} says how beliefs are in fact formed, sustained and revised, by what processes, and with what frequency of truth in what environments. A \emph{normative epistemic claim} says what a subject is justified in believing, what she ought to believe, or whether a belief is knowledge.
\end{definition}

Knowledge is a normative status in the sense of Definition~\ref{d:desc}: to say that a belief is knowledge is to credit it with a standing that true belief may lack. We take knowledge to be safe, justified true belief (EPI-01), and justification to be internalist, while knowledge is externalist (EPI-06).

\begin{principle}[Epistemic ought implies can]\label{p:can}
If a subject ought epistemically to believe, suspend or reason in a certain way, then she is able to do so.
\end{principle}

Principle~\ref{p:can} is the epistemic counterpart of the familiar ethical principle. It has to be stated with care: a subject can be bound by norms that she fails to satisfy, and what is required is the capacity, under ordinary conditions, not success. It excludes norms that require computation beyond any finite agent, such as deductive closure of belief or probabilistic coherence over infinitely many propositions, from being norms that bind human believers as such, though they may remain ideals by which performance is measured.

\begin{principle}[Safety is modal and empirical]\label{p:safe}
Whether a belief formed on a given basis is safe depends on whether, in close worlds, beliefs formed on that basis are true; this depends on the reliability of the basis in the subject's environment.
\end{principle}

Principle~\ref{p:safe} follows from the definition of safety in EPI-01: a belief is safe if, in every close world in which the subject believes the same proposition on the same basis, it is true. Whether this holds for perceptual beliefs in poor light, for memory of faces, or for testimony from experts in a given field depends on facts that are discovered, when they are discovered, by experiment.

\begin{principle}[Bridge principles]\label{p:bridge}
Any argument from descriptive premises to a normative epistemic conclusion has, explicitly or implicitly, a premise linking a descriptive property of beliefs or processes to a normative status. A \emph{bridge principle} is such a premise.
\end{principle}

Examples of bridge principles are \enquote{a belief that is safe and justified is knowledge}, \enquote{a seeming that $p$ confers prima facie justification on believing $p$} (EPI-19) and \enquote{a process whose outputs are mostly true in the subject's environment confers justification} \citep{Goldman1979}. Principle~\ref{p:bridge} does not say that normative facts float free of descriptive ones; it is compatible with the supervenience of the normative on the descriptive, and indeed bridge principles are what such supervenience consists in. It says that the inference requires them.

\section{Why epistemology needs psychology}

\begin{theorem}\label{t:need}
Epistemology cannot settle which beliefs of actual human subjects are knowledge, or which norms of reasoning bind them, without psychological premises.
\end{theorem}

\begin{proof}
There are two routes, one for each part of the claim.

\emph{Knowledge.} By the analysis of EPI-01, a belief is knowledge only if it is safe, and by Principle~\ref{p:safe} whether it is safe depends on the reliability of its basis in the subject's environment. Whether human face recognition across ethnic groups, eyewitness memory after leading questions, or confidence-based judgment in a given domain is reliable is an empirical matter. So whether the beliefs formed by these means are knowledge cannot be settled without empirical premises about human cognition, which are psychological premises.

\emph{Norms.} By Principle~\ref{p:can}, a norm binds human believers only if they can comply with it. What human believers can do, with what effort and under what conditions, is a matter for psychology: the limits of working memory, the availability of heuristics and the conditions under which they can be overridden are empirical findings \citep{TverskyKahneman1974}. So which candidate norms bind human believers cannot be settled without psychological premises.
\end{proof}

The second route has substantive consequences. A norm requiring believers to update by Bayesian conditionalization on every item of evidence binds no human believer as stated; what binds is a norm requiring them to approximate it by procedures they can carry out, and which procedures approximate it best for creatures like us is an empirical question. This is the core of truth in the programme of \citet{BishopTrout2005}: the finding that simple linear rules outperform expert clinical judgment in many domains is a finding about which reasoning strategies are, for us, the better ones.

\section{Why the normative questions remain distinct}

\begin{theorem}\label{t:distinct}
No set of purely descriptive claims about human cognition entails a normative epistemic claim; the bridge principles required are not findings of psychology.
\end{theorem}

\begin{proof}
By Principle~\ref{p:bridge}, an argument from descriptive premises to a normative conclusion requires a bridge principle. A bridge principle links a descriptive property, such as reliability, safety or the presence of a seeming, to a normative status, such as justification or knowledge. Psychology, as the empirical study of cognition, establishes which processes have which descriptive properties: how often eyewitness memory is accurate under given conditions, which heuristics people use. It does not establish which descriptive properties confer justification, since its methods (experiment, observation, modelling) bear on what processes are like and not on what standing their outputs have. Two theorists who agree on every psychological finding can disagree on whether reliability or the presence of an undefeated seeming is what confers justification, as externalists and internalists about justification do; no further experiment would settle their disagreement.
\end{proof}

\begin{corollary}\label{c:both}
Epistemology is continuous with psychology, since its conclusions about human knowledge and about the norms binding human believers depend on psychological premises; and its normative questions are distinct from psychological ones, since its conclusions depend also on bridge principles that psychology does not supply.
\end{corollary}

Quine's later formulation, on which normative epistemology is the technology of truth-seeking, fits Corollary~\ref{c:both} better than his earlier one. A technology presupposes an end: engineering asks how to build bridges that stand, and the value of bridges that stand is not a finding of engineering. To say that the end of inquiry is truth, or truth reached non-accidentally, is already to take a normative position, and the choice between ends of this kind is what epistemology adds to psychology.

\section{Objections and replies}

\begin{objection}[Bridge principles are empirical]\label{o:empirical}
A naturalist may reply that bridge principles are themselves empirical hypotheses, to be accepted or rejected according to their fit with our practice of epistemic evaluation, which is a natural phenomenon \citep{Kornblith2002}.
\end{objection}

\begin{reply}
If the evidence for a bridge principle is that we in fact evaluate beliefs in accordance with it, the principle so confirmed is descriptive: it says how we evaluate, not that the evaluation is correct. If the evidence is that beliefs meeting the principle are more often true, the principle presupposes that truth is the relevant end, which is a normative premise. Our practice of evaluation is evidence for bridge principles, as moral intuitions are evidence for moral principles, but it is evidence because it is taken to track something it may fail to track.
\end{reply}

\begin{objection}[Ought implies can is too weak]\label{o:weak}
Principle~\ref{p:can} constrains which norms bind, but only negatively. Psychology then contributes nothing to the content of epistemic norms, and the continuity claimed in Theorem~\ref{t:need} is minimal.
\end{objection}

\begin{reply}
The contribution is not merely negative. Once ideal norms are excluded as binding, the binding norms are those that best serve the epistemic ends given our capacities, and which strategies serve those ends best is a comparative empirical question, as the case of statistical prediction rules shows. Psychology thus determines, given the ends, the content of the norms that bind us. Only the ends come from elsewhere.
\end{reply}

\begin{objection}[Kim's dilemma]\label{o:kim}
If justification is normative, as we hold, then the naturalist cannot help herself to it in describing cognition, and if it is descriptive, the normativity is lost. There is no stable middle position \citep{Kim1988}.
\end{objection}

\begin{reply}
The middle position is the one defended above: psychology describes cognition without using the concept of justification, and epistemology combines its descriptions with bridge principles to reach normative conclusions. Neither discipline needs the other's concepts in its own premises; epistemology needs psychology's conclusions. This is how applied ethics uses facts about consequences, and it no more makes epistemology a part of psychology than it makes ethics a part of economics.
\end{reply}

\section{Conclusion}

Epistemology cannot do without psychology. Whether human beliefs are knowledge depends on the reliability of the processes that form them, and which norms bind human believers depends on what they are able to do; both are empirical questions. But the normative questions of epistemology are not answered by psychology, since every passage from facts about cognition to verdicts about justification and knowledge passes through bridge principles that are not psychological findings. Quine was right that epistemology belongs within science, in the sense that it takes science's results as premises, and Kim was right that it does not become a description of how we come to believe.

\begin{thebibliography}{99}\small
\bibitem[Bishop and Trout(2005)]{BishopTrout2005} Bishop, M.~A. and Trout, J.~D. (2005). \emph{Epistemology and the Psychology of Human Judgment}. Oxford University Press.
\bibitem[Goldman(1979)]{Goldman1979} Goldman, A.~I. (1979). What is justified belief? In G.~S. Pappas (Ed.), \emph{Justification and Knowledge} (pp.~1--23). Reidel.
\bibitem[Goldman(1986)]{Goldman1986} Goldman, A.~I. (1986). \emph{Epistemology and Cognition}. Harvard University Press.
\bibitem[Kim(1988)]{Kim1988} Kim, J. (1988). What is \enquote{naturalized epistemology}? \emph{Philosophical Perspectives}, 2, 381--405.
\bibitem[Kornblith(2002)]{Kornblith2002} Kornblith, H. (2002). \emph{Knowledge and Its Place in Nature}. Oxford University Press.
\bibitem[Quine(1969)]{Quine1969} Quine, W.~V. (1969). Epistemology naturalized. In \emph{Ontological Relativity and Other Essays} (pp.~69--90). Columbia University Press.
\bibitem[Quine(1986)]{Quine1986} Quine, W.~V. (1986). Reply to Morton White. In L.~E. Hahn and P.~A. Schilpp (Eds.), \emph{The Philosophy of W. V. Quine}. Open Court.
\bibitem[Stich(1990)]{Stich1990} Stich, S. (1990). \emph{The Fragmentation of Reason}. MIT Press.
\bibitem[Tversky and Kahneman(1974)]{TverskyKahneman1974} Tversky, A. and Kahneman, D. (1974). Judgment under uncertainty: Heuristics and biases. \emph{Science}, 185(4157), 1124--1131.
\end{thebibliography}

\end{document}
